% Ban thu: chi giu cac cau co trong CD2

%@ id=TT26.thaithinh.003 | trang=1 | so=2.1 | de=II | vitri=CD2.4,C3.9.5 | muc=8 | tt=0
\begin{bt}
Cho hai biểu thức $A = \dfrac{x+\sqrt{x}+1}{\sqrt{x}+2}$ và $B = \dfrac{1}{\sqrt{x}-2} + \dfrac{x-\sqrt{x}-4}{x-2\sqrt{x}}$ với $x > 0; x \neq 4$.
\begin{cacy}(1)
\item Tính giá trị của biểu thức $A$ khi $x = 9$.
\item Chứng minh $B = \dfrac{\sqrt{x}+2}{\sqrt{x}}$.
\item Cho $P = AB$. So sánh $P$ với 3.
\end{cacy}
\loigiai{
\begin{cacy}(1)
\item Thay $x = 9$ (TM ĐK) vào biểu thức $A$, ta có $A = \dfrac{9+\sqrt{9}+1}{\sqrt{9}+2} = \dfrac{13}{5}$.
\item \[
B = \dfrac{\sqrt{x}}{\sqrt{x}-2} + \dfrac{x-\sqrt{x}-4}{\sqrt{x}(\sqrt{x}-2)} = \dfrac{x-4}{\sqrt{x}(\sqrt{x}-2)}
\]
\[
B = \dfrac{(\sqrt{x}+2)(\sqrt{x}-2)}{\sqrt{x}(\sqrt{x}-2)} = \dfrac{\sqrt{x}+2}{\sqrt{x}}
\]
\item $P = AB = \dfrac{x+\sqrt{x}+1}{\sqrt{x}}$. Xét hiệu $P - 3 = \dfrac{x+\sqrt{x}+1}{\sqrt{x}} - 3 = \dfrac{(\sqrt{x}-1)^2}{\sqrt{x}}$.

Với ĐKXĐ: $x > 0; x \neq 4$ thì $\dfrac{(\sqrt{x}-1)^2}{\sqrt{x}} \ge 0$.

Vậy $P \ge 3$.
\end{cacy}
}
\end{bt}
